\chapter{Predictive Parsing and Grammar Design}
\label{ch:predictive-grammar}

An identifier can begin an assignment, a call, or a bare expression.
A parser that selects the first plausible production can discover its mistake only after consuming a shared prefix.
One response is to retain alternatives and backtrack.
Another is to design a grammar whose alternatives can be selected before any of them runs.
This chapter makes the second response mechanical: compute what may begin a production, what may follow it, and whether those sets determine a unique choice.
The result is both an implementation recipe and a diagnostic for the grammar author.

\section{A lookahead decision is a proof obligation}
\label{sec:ch50-lookahead}

We use token kinds rather than source characters in this chapter.
Let the grammar describe a bracketed list of identifiers:
\begin{align*}
 S&\to [\ I\ ],\\
 I&\to \mathit{id}\ T\mid\epsilon,\\
 T&\to ,\ \mathit{id}\ T\mid\epsilon.
\end{align*}
The symbols \code{[}, \code{]}, and \code{,} are terminals; $S,I,T$ are nonterminals.
The external end marker $\$$ is not a grammar token that the lexer can emit in the middle of a list.
The language contains \code{[]}, \code{[id]}, and \code{[id,id]}, but excludes \code{[id,]}.

At $I$, an identifier selects the nonempty alternative.
A closing bracket selects the empty alternative because it belongs to the caller.
An unrelated token must not select epsilon merely because the parser cannot begin the nonempty alternative.
Doing so loses the distinction between a legal empty list and an error.
Figure~\ref{fig:ch50-lookahead} draws the decision at that boundary.

\begin{figure}[htbp]
\centering
\begin{tikzpicture}[node distance=15mm and 20mm]
\node[basic block] (i) {parse I};
\node[basic block,below left=of i] (id) {id T};
\node[basic block,below=of i] (empty) {epsilon};
\node[basic block,below right=of i] (error) {error};
\draw[cfg edge] (i)--node[left,font=\small] {\code{id}}(id);
\draw[cfg edge] (i)--node[right,font=\small] {\code{]}}(empty);
\draw[back edge] (i)--node[right,font=\small] {other}(error);
\end{tikzpicture}
\caption{Epsilon is selected by the caller's permitted follower, not by a catch-all failure branch.}
\label{fig:ch50-lookahead}
\end{figure}

An \term{LL(1) grammar} permits left-to-right scanning and a leftmost derivation with one token of lookahead, under the usual predictive table construction.
The qualification concerns the grammar, not only its language: a language may have both a convenient LL(1) grammar and another grammar with conflicts.
Increasing lookahead is also not equivalent to ignoring conflicts and hoping the first rule works~\cite{grune-2008-parsing}.

\section{Compute nullable and FIRST monotonically}
\label{sec:ch50-first}

A nonterminal is \term{nullable} when it derives the empty string.
The set $\operatorname{FIRST}(X)$ contains terminals that can begin a nonempty string derived from $X$.
We keep nullability separate, so epsilon is never accidentally used as a token-table column.
For a sequence, inspect its symbols from left to right, adding their FIRST sets until reaching a nonnullable symbol.
The empty sequence is nullable and has an empty FIRST set.

Algorithm~\ref{alg:ch50-first} handles both properties in one monotone iteration.
For a terminal $a$, initialize $\operatorname{FIRST}(a)=\{a\}$ and mark it nonnullable.
For nonterminals, initialize both properties to empty or false.
The algorithm must revisit a production when a symbol later becomes nullable; a single textual pass is insufficient.

\begin{algorithm}[htbp]
\caption{Compute nullable flags and terminal FIRST sets to a fixed point.}
\label{alg:ch50-first}
\begin{pseudocode}
\Require finite grammar with defined symbols
\Ensure nullable and FIRST for every nonterminal, excluding epsilon from FIRST
\State initialize terminal facts; initialize nonterminal facts to empty/false
\Repeat
  \State $changed\gets false$
  \ForAll{productions $A\to X_1\cdots X_k$}
    \State $all\gets true$
    \For{$j=1,\ldots,k$}
      \State union $\operatorname{FIRST}(X_j)$ into $\operatorname{FIRST}(A)$
      \If{$X_j$ is not nullable}
        \State $all\gets false$; \Break
      \EndIf
    \EndFor
    \If{$all$}
      \State mark $A$ nullable
    \EndIf
    \State record whether any set or flag changed
  \EndFor
\Until{nothing changed}
\end{pseudocode}
\end{algorithm}

There are at most $N(t+1)$ facts for $N$ nonterminals and $t$ terminal kinds.
Every changing sweep adds a fact.
If $P$ is the number of productions, $L$ the total number of RHS symbol occurrences, and $W$ the machine-word count of a terminal bitset, a direct implementation costs at most $O((N(t+1)+1)(P+LW))$ word operations.
The $P$ term covers empty productions too.
Dependency worklists can improve this coarse full-sweep bound.

The proof follows derivation trees.
Each inserted fact has a finite justification using productions and facts already known.
Conversely, induction on the height of a finite derivation eventually inserts its nullability or first-token fact.
Recursive rules with no terminating derivation do not acquire facts simply by referring to themselves.
For $A\to A$, the result is neither nullable nor able to start with any terminal.

\section{FOLLOW carries the caller's obligations}
\label{sec:ch50-follow}

The set $\operatorname{FOLLOW}(A)$ contains terminal kinds, or the end marker, that may appear immediately after an occurrence of $A$ in a derivation from the start symbol.
For the standard table construction, unreachable productions can be removed first; otherwise the usual equations may include conservative extra context from unreachable rules.
Seed $\operatorname{FOLLOW}(S)$ with the end marker.
An occurrence $B\to\alpha A\beta$ contributes $\operatorname{FIRST}(\beta)$ to $\operatorname{FOLLOW}(A)$.
If $\beta$ is nullable, it also contributes $\operatorname{FOLLOW}(B)$.
Algorithm~\ref{alg:ch50-follow} makes both contributions explicit.

\begin{algorithm}[htbp]
\caption{Propagate followers through nullable suffixes.}
\label{alg:ch50-follow}
\begin{pseudocode}
\Require reachable grammar, completed nullable/FIRST facts, start $S$
\Ensure FOLLOW sets including the external end marker
\State initialize all FOLLOW sets empty; insert end marker in $\operatorname{FOLLOW}(S)$
\Repeat
  \ForAll{nonterminal occurrences $B\to\alpha A\beta$}
    \State union $\operatorname{FIRST}(\beta)$ into $\operatorname{FOLLOW}(A)$
    \If{$\beta$ is nullable}
      \State union $\operatorname{FOLLOW}(B)$ into $\operatorname{FOLLOW}(A)$
    \EndIf
  \EndFor
\Until{no set changed}
\end{pseudocode}
\end{algorithm}

Precompute suffix FIRST sets and nullable flags by right-to-left scans of each RHS.
Then a sweep needs $O(LW)$ word operations rather than repeatedly rescanning long suffixes.
With at most $N(t+1)$ follower memberships, the analogous full-sweep bound is $O((N(t+1)+1)LW)$ after preprocessing.
Table~\ref{tab:ch50-sets} gives the exact facts for our list grammar.
The closing bracket flows from $S$ to $I$, then through the nullable suffix of $I\to\mathit{id}\ T$ to $T$.
The comma belongs to FIRST of $T$, not FOLLOW of $I$.

\begin{table}[htbp]
\caption{The list grammar separates nonempty beginnings from legal empty returns.}
\label{tab:ch50-sets}
\centering
\begin{tabular}{llll}
\toprule
Nonterminal & Nullable & FIRST & FOLLOW \\
\midrule
$S$ & no & $\{[\}$ & $\{\$\}$ \\
$I$ & yes & $\{\mathit{id}\}$ & $\{]\}$ \\
$T$ & yes & $\{,\}$ & $\{]\}$ \\
\bottomrule
\end{tabular}
\end{table}

\section{Construct the table and expose conflicts}
\label{sec:ch50-table}

For each production $A\to\alpha$, its selection set is
\[
 \operatorname{SELECT}(A\to\alpha)=\operatorname{FIRST}(\alpha)
 \cup
 \begin{cases}
 \operatorname{FOLLOW}(A),&\alpha\text{ nullable},\\
 \emptyset,&\text{otherwise}.
 \end{cases}
\]
Insert the production into every table cell $(A,a)$ named by that set.
A cell with two distinct productions is a conflict to report, not an invitation to overwrite the earlier entry.
Table~\ref{tab:ch50-table} lists all populated cells; all other cells are errors.

\begin{table}[htbp]
\caption{Every populated predictive-table cell contains exactly one production.}
\label{tab:ch50-table}
\centering
\begin{tabular}{lll}
\toprule
Nonterminal & Lookahead & Production \\
\midrule
$S$ & \code{[} & $S\to[\ I\ ]$ \\
$I$ & identifier & $I\to\mathit{id}\ T$ \\
$I$ & \code{]} & $I\to\epsilon$ \\
$T$ & \code{,} & $T\to,\ \mathit{id}\ T$ \\
$T$ & \code{]} & $T\to\epsilon$ \\
\bottomrule
\end{tabular}
\end{table}

Two different failure modes matter.
Alternatives $A\to\mathit{id}\ =\ E\mid\mathit{id}\ (\ Args\ )$ share FIRST and need factoring or more lookahead.
In $A\to B\mid\epsilon$ with $B$ starting in a token also in FOLLOW of $A$, the conflict involves the empty return path.
Factoring a common nonempty prefix does not automatically remove this second kind of conflict.

Left factoring rewrites the first example as $A\to\mathit{id}\ R$ with $R\to=\ E\mid(\ Args\ )$.
The decision moves to the point after the identifier.
Figure~\ref{fig:ch50-factoring} makes that relocation visible.
It preserves recognition, but the semantic action building the assignment or call must now obtain the identifier from the outer context.
Grammar rewriting without corresponding tree-construction changes can preserve accepted strings and still build wrong ASTs.

\begin{figure}[htbp]
\centering
\begin{tikzpicture}[x=24mm,y=12mm]
\node[basic block] (old) at (0,1) {A};
\node[basic block] (o1) at (2,2) {id = E};
\node[basic block] (o2) at (2,0) {id (Args)};
\draw[cfg edge] (old)--(o1); \draw[cfg edge] (old)--(o2);
\node[basic block] (new) at (0,-2) {A: id R};
\node[basic block] (n1) at (2,-1) {R: = E};
\node[basic block] (n2) at (2,-3) {R: (Args)};
\draw[cfg edge] (new)--node[above,font=\small] {lookahead =}(n1);
\draw[cfg edge] (new)--node[below,font=\small] {lookahead (}(n2);
\node[font=\small,anchor=west] at (3,1) {same first token};
\node[font=\small,anchor=west] at (3,-2) {decision after id};
\end{tikzpicture}
\caption{Factoring delays a choice until the distinguishing token is visible, while the shared identifier is consumed once.}
\label{fig:ch50-factoring}
\end{figure}

\section{Drive an explicit predictive stack}
\label{sec:ch50-driver}

A recursive-descent implementation can encode each table row as a switch.
An explicit stack instead stores the unprocessed sentential suffix.
Algorithm~\ref{alg:ch50-driver} gives a recognizer; AST construction adds reduction markers carrying children and source spans.
The grammar must be well formed and have a conflict-free predictive table, without useless nonprogress cycles.

\begin{algorithm}[htbp]
\caption{Recognize with a conflict-free predictive table.}
\label{alg:ch50-driver}
\begin{pseudocode}
\Require token sequence with one end marker, start $S$, predictive table $M$
\Ensure accept exactly when the entire sequence derives from $S$
\State stack $\gets [S,end]$, with top at the left; cursor $i\gets0$
\While{stack is nonempty}
  \State $X\gets$ pop stack; $a\gets$ lookahead at $i$
  \If{$X$ is a terminal or end marker}
    \If{$X\neq a$} \State \Return syntax error \EndIf
    \State $i\gets i+1$
  \Else
    \If{$M[X,a]$ is empty} \State \Return syntax error \EndIf
    \State push the selected RHS in reverse order, omitting epsilon
  \EndIf
\EndWhile
\State \Return accept
\end{pseudocode}
\end{algorithm}

On \code{[id,id]}, the stack begins $[S,\$]$.
After expanding $S$ and matching the opening bracket it is $[I,],\$]$.
Selecting and matching the first identifier leaves $[T,],\$]$.
The comma selects $T\to,\ \mathit{id}\ T$; after those two terminals the same stack shape recurs with a closing-bracket lookahead.
The epsilon production pops $T$, then the bracket and end marker match.
The expansion sequence is $S$, nonempty $I$, nonempty $T$, empty $T$.
The cursor never moves backward.

For a fixed suitable grammar, predictive parsing takes linear time in the token count: each expansion contributes bounded grammar work, and nonprogress expansion chains have a grammar-dependent bound.
When grammar size is also variable, state the bound as token matches plus the total length of expanded right-hand sides, rather than hiding large productions in a constant.
AST storage is proportional to the nodes actually constructed; the explicit pending-symbol stack can still grow with input nesting.

\section{Associativity lives in the AST construction}
\label{sec:ch50-associativity}

The grammar $E\to E-\mathit{num}\mid\mathit{num}$ is directly left recursive.
A naive top-down function calls itself before consuming input.
The language-equivalent form $E\to\mathit{num}\ R$, $R\to-\mathit{num}\ R\mid\epsilon$ terminates predictively, but its right-recursive concrete tree does not prescribe a right-associative AST.
Build an initial numeric node and fold each encountered subtraction into a new left child/root pair.
For \code{20-6-3}, the left fold computes $(20-6)-3=11$.
A right-associated AST computes $20-(6-3)=17$.
Figure~\ref{fig:ch50-assoc} shows why recognition alone cannot test the transformation.

\begin{figure}[htbp]
\centering
\begin{tikzpicture}[x=11mm,y=11mm]
\node[basic block] (l) at (0,0) {Sub};
\node[basic block] (ll) at (-1.5,-1.5) {Sub};
\node (a) at (-2.5,-3) {20}; \node (b) at (-0.5,-3) {6}; \node (c) at (1.5,-1.5) {3};
\draw[cfg edge] (l)--(ll); \draw[cfg edge] (ll)--(a); \draw[cfg edge] (ll)--(b); \draw[cfg edge] (l)--(c);
\node[basic block] (r) at (6,0) {Sub};
\node (d) at (4.5,-1.5) {20}; \node[basic block] (rr) at (7.5,-1.5) {Sub};
\node (e) at (6.5,-3) {6}; \node (f) at (8.5,-3) {3};
\draw[cfg edge] (r)--(d); \draw[cfg edge] (r)--(rr); \draw[cfg edge] (rr)--(e); \draw[cfg edge] (rr)--(f);
\node at (0,-4) {left fold: $11$}; \node at (6,-4) {right grouping: $17$};
\end{tikzpicture}
\caption{The same token sequence permits different results unless AST construction preserves the intended associativity.}
\label{fig:ch50-assoc}
\end{figure}

Precedence layers similarly encode which subparser may consume which operators.
Pratt parsing and precedence climbing express this through binding powers rather than separate grammar functions.
Either approach needs explicit policies for equal precedence, prefix operators, postfix operators, and nonassociative comparisons.
Test tree shapes, not just whether a string was accepted.

\section{Recovery must preserve progress and report damage}
\label{sec:ch50-recovery}

On \code{[id,]}, the comma has already committed the parser to another identifier.
The mismatch at the bracket is therefore an expected-identifier error, not permission to return from the tail.
For a compiler continuing after errors, one possible policy inserts an error-node identifier before that bracket and lets the caller consume the bracket normally.
The recovered tree is marked invalid for code generation.
It is not a successful parse of the original grammar.

\begin{algorithm}[htbp]
\caption{Synchronize a failed nonterminal without consuming a caller's delimiter.}
\label{alg:ch50-recover}
\begin{pseudocode}
\Require failed nonterminal $A$, token cursor, synchronization set containing end
\Ensure invalid placeholder and progress by popping $A$ or consuming input
\State emit one diagnostic at the current token
\While{lookahead is not in the synchronization set}
  \State consume one token
\EndWhile
\State discard the pending expansion of $A$
\State return an error node without consuming the synchronization token
\end{pseudocode}
\end{algorithm}

Algorithm~\ref{alg:ch50-recover} terminates because the end marker belongs to its stop set.
The caller must actually remove the failed stack frame; retrying the same expansion on the same delimiter would loop.
FOLLOW sets provide candidate synchronization tokens, but nesting and language-specific statement boundaries often require a more selective policy.
Diagnostic suppression for cascades is an interface decision and must not silently turn a damaged tree into a valid compilation unit.

\section{Exercises}
\label{sec:ch50-exercises}
\begin{enumerate}
\item\label{ex:ch50-three} Trace \code{[id,id,id]} with Algorithm~\ref{alg:ch50-driver}. Count production expansions and terminal matches, including end.
\item\label{ex:ch50-conflict} For $S\to A\ \mathit{id}$ and $A\to\mathit{id}\mid\epsilon$, compute SELECT of both $A$ productions and identify the conflicting cell.
\item\label{ex:ch50-fold} Compute left and right groupings of \code{30-8-5}; describe the accumulator update that selects the left grouping.
\item\label{ex:ch50-recovery} On \code{[id,,id]}, synchronize at comma, bracket, or end. Explain why popping the failed identifier obligation is necessary before leaving the comma unconsumed.
\end{enumerate}

\section{Further reading}
\label{sec:ch50-reading}
Grune and Jacobs compare grammar transformations and predictive methods, emphasizing the difference between a grammatical language and a parsing strategy~\cite{grune-2008-parsing}.
Hopcroft, Motwani, and Ullman supply the context-free grammar and pushdown-machine background~\cite{hopcroft-2007-automata}.
